% Pursuing Stacks (cotes 134-2 to 134-8): archive cote 134-4, site page 51 (PDF page 52), typescript page 139.
% typescript. Draft compared: ps3.tex:2573-2650.
% AI-transcribed by Claude Opus 5.5 (claude-opus-5-5, adaptive thinking, effort high), Message Batches API, batch ps2 (msgbatch_01NVwr1Yuo8qJf5rykeGBXPK), 2026-11-12; not yet checked by eye.
\scanpage{134-4/51}
\mspage{139}

beginning  -  namely an\xout{xxx} investigation of modelizers and,  more specifi-%
cally,  the gradual unravelling of a notion of ``canonical modelizer''.  I
certainly feel like carrying this to the end at once,  without any digres-%
sions except when felt relevant for the main focus at present.  As for
choosing precedence between c) and d),  it is still time to decide,  when
we're through with a) and b) !

5\over{\unsure{6}}{7}.  \uline{Relation of homotopy structures to} 0-\uline{connectedness and to} $\hand{\pi}_0$ .

Here we're resuming the reflection started in par.51 F) (page 122).
All we did there was introduce some conditions on a category M,  namely
a) to e)  (page 123),  and introduce the functor
\[ \hand{\pi}_0 :\ M \hand{\longrightarrow} (\mathrm{Sets}) \qquad . \]
As before,  we'll assume now M satisfies conditions a) to d)  (the first
three,\xout{xxx} or rather b) and c),  are all that is needed for defining the
functor $\hand{\pi}_0$) ,  and will not assume e),  or the equivalent condition e')
of $\hand{\pi}_0$ commuting to \xout{\unsure{pro}} finite products,  unless explicitly specified.

We suppose now,  moreover,  M endowed with a homotopy structure\xout{s} h.
We'll say that h is $\hand{\pi}_0$-\uline{admissible},  or \inserted{simply} 0-\uline{admissible},  if \xout{th} h is compatib-%
le with the functor $\hand{\pi}_0$ ,  which can be expressed by either one of the
following two equivalent conditions\xout{x} (cf. page 130):

(i)  \xout{t}f $\underset{h}{\hand{\sim}}$ g  implies $\hand{\pi}_0$(f) = $\hand{\pi}_0$(g)  ,

for any two maps \xout{xx} f,g in M,  or

(ii)  f $\hand{\in} W_h$ (i.e. f a homotopism)  implies $\hand{\pi}_0$(f) bijective.

Another equivalent formulation is that the functor $\hand{\pi}_0$ factors through
the quotient category $M_h$ of homotopy types
\[ M \hand{\longrightarrow} M_h \hand{\longrightarrow} (\mathrm{Sets}) \qquad , \]
we'll still denote by $\hand{\pi}_0$ the functor $M_h \hand{\longrightarrow}$ (Sets)  obtained.

If h is 0-admissible,  it satisfies  (i') and (ii') below:

(i')  \xout{f}For any weak homotopy interval $(I,\hand{\delta}_0,\hand{\delta}_1)\hand{\in}\hand{\Sigma}_h$ , $\hand{\pi}_0(\hand{\delta}_0)$ = $\hand{\pi}_0(\hand{\delta}_1)$.

As $\hand{\pi}_0(e_M)$  is a one-point set by the assumption d) on M,  for any section
$\hand{\delta}$ of an object I , $\hand{\pi}_0(\hand{\delta})$ may be described as just an element of $\hand{\pi}_0$(I),
which is the unique connected component of I through which factors the
given section;  thus (i') can be expressed by saying that for any weak
homotopy interval, the two ``endpoints'' belong to the same connected compo-%
nent of I, \xout{wxx} a natural condition indeed !  It is automatically satis-%

\margin{51}
